En un problema de programación lineal de maximización, dada una restricción de tipo menor o igual, ¿qué se puede afirmar sobre el signo del precio sombra de dicha restricción correspondiente una solución básica y óptima del problema?
\vspace{1cm}

\noindent
\textbf{Solución}

\vspace{0.5cm}

Cualquier solución básica óptima cumple que $V^B \leq 0$. Por otro lado, el precio sombra de una restricción menor o igual es igual al criterio del simplex correspondiente a la variable de holgura de dicha restricción cambiado de signo, es decir, $\pi^B_i = -V^B_{h_i}$. Por lo tanto, $\pi^B_i \geq 0$, el precio sombra es no negativo en la solución óptima.
